\section{2. Cơ sở lý thuyết về mạng WLAN}

\subsection{2.1 Khái niệm}


\subsection{2.2 Các thành phần của mạng WLAN}
Mạng WLAN được cấu thành từ nhiều thiết bị và thành phần khác nhau, đóng vai trò quan trọng trong việc truyền tải và xử lý tín hiệu không dây:
\begin{itemize}
    \item \textbf{Access Point (AP) / Wireless Router:}
    \item \textbf{Wireless Network Interface Controller (WNIC):} 
    \item \textbf{Thiết bị đầu cuối (End Devices):} 
    \item \textbf{Môi trường truyền dẫn:} 
\end{itemize}

\subsection{2.3 Mô hình triển khai và Giải pháp kết nối}

\subsubsection{2.3.1 Mô hình Infrastructure (Có Access Point/Router trung tâm)}

\begin{itemize}
    \item \textbf{Ưu điểm:} 
    \item \textbf{Nhược điểm:} 
\end{itemize}

\subsubsection{2.3.2 Mô hình Ad-hoc (Mạng ngang hàng)}

\begin{itemize}
    \item \textbf{Ưu điểm:} 
    \item \textbf{Nhược điểm:} 
\end{itemize}

\subsubsection{2.3.3 Giải pháp kết nối cho thiết bị đầu cuối}
Dưới đây là các giải pháp kết nối mạng và cấu hình cho 3 loại thiết bị phổ biến:
\begin{itemize}
    \item \textbf{Laptop:} 
    \item \textbf{Desktop (PC):}
    \item \textbf{Smartphone:} 
\end{itemize}